\chapter{系统测试与结果分析}

\section{测试目标}

系统测试的目标，是验证布控球自适应 PTZ 跟踪原型是否能够完成预期功能，并分析当前实现中的功能完整性、实时性和控制表现。测试内容围绕设备接入、视频流获取、人体检测、ReID 目标记忆、PTZ 自动跟踪、交互模块和资源释放展开。由于本文系统依赖真实布控球设备和厂商 SDK 接口，测试重点主要放在工程闭环是否稳定、关键模块是否能够按预期协同，以及系统运行状态是否符合设计要求。

测试过程采用由局部到整体的思路。前期先验证设备登录、实时预览和解码链路是否正常，确保系统具有稳定的视频输入；随后验证 YOLO 检测、目标编号和目标选择逻辑，确认视觉分析模块能够输出可用于控制的数据；之后再验证 PTZ 自动跟踪和界面交互，观察系统是否能够把检测结果转化为真实云台动作。通过这种分层测试方式，系统出现异常时能够更快定位问题来源。

由于本文系统属于真实设备原型，测试过程并不完全等同于离线算法评测。离线算法评测通常可以在固定数据集上重复运行，并用统一指标比较模型优劣；而布控球跟踪系统还受到网络状态、设备机械响应、安装姿态和现场光照等因素影响。因此，本文测试更关注系统链路是否能够稳定贯通，以及各模块之间的数据是否能够按预期传递。对于暂时无法完全量化的部分，则通过日志、界面反馈和人工观察进行辅助判断。

\section{测试环境}

当前系统的测试环境见表\ref{tab:test-env}。

\begin{table}[H]
    \centering
    \caption{系统测试环境}
    \label{tab:test-env}
    \begin{tabularx}{\textwidth}{c|X}
        \hline
        项目 & 配置 \\
        \hline
        验证平台 & 原型验证主机 \\
        \hline
        开发语言 & Python \\
        \hline
        设备条件 & 海康威视布控球，与验证主机保持稳定网络连接 \\
        \hline
        设备接口 & HCNetSDK、PlayCtrl \\
        \hline
        检测模型 & YOLOv8n，输入尺寸 384，置信度阈值 0.4 \\
        \hline
        视觉计算库 & ultralytics、opencv-python、Pillow、NumPy \\
        \hline
    \end{tabularx}
\end{table}

测试时，验证主机与布控球保持稳定网络连接，以减少网络波动对视频预览和控制指令的影响。测试前完成设备连通性、接口初始化和模型加载能力检查，保证视频采集链路和检测链路处于可用状态。由于不同硬件配置会影响 YOLO 推理速度，本文测试结论主要用于说明原型系统的可行性，不作为严格硬件基准。

\section{功能测试}

\subsection{设备登录与视频预览测试}

设备接入测试用于验证 SDK 初始化、设备登录和实时预览流程。测试过程中重点观察设备接口是否能够返回有效登录状态，预览链路是否能够持续输出码流，以及解码后的图像帧是否能够被系统稳定读取。根据中期检查结果，系统已经实现从 SDK 初始化、设备登录、启动预览、帧回调到解码显示的完整流程，实时视频流能够稳定获取，延迟控制在 200 ms 以内。

在测试过程中，设备接入状态主要与网络连通性、接口初始化、播放库端口分配和码流回调链路有关。该测试不仅验证功能是否可用，也帮助明确真实设备接入时影响系统稳定性的关键环节。

设备登录与视频预览测试还承担了后续模块测试的前置作用。视频链路保持稳定后，检测结果、目标选择和云台控制才具有可靠输入基础；预览句柄和播放库端口按照流程释放后，也能够保证系统再次启动时保持清晰的初始化状态。因此，该测试不仅是功能验证，也是保障后续模块顺利联调的基础步骤。

\subsection{YOLO 人体检测测试}

检测测试用于验证系统能否识别画面中的人体目标。在检测模块启用状态下，系统应输出人体边界框、目标 ID、置信度和中心点等信息。根据中期检查材料，系统选用 YOLOv8n 轻量级模型，在当前配置下能够实现约 30 fps 的实时检测效果。从这一结果看，轻量模型基本能够满足原型系统的实时闭环需求。

检测结果反馈还用于观察检测框稳定性和置信度变化。当目标距离较近、光照较稳定且遮挡较少时，YOLOv8n 能够持续输出人体框；当目标距离较远、背景复杂或画面运动较快时，检测框会随画面变化产生短时波动。针对这一现象，系统没有直接把单帧检测结果作为最终控制输入，而是结合 ReID 记忆、目标选择和平滑控制机制共同使用，使控制输入更加稳健。

在检测测试中，本文重点关注人体目标是否能够被持续识别，以及检测框中心点是否能作为控制输入使用。对于布控球云台控制而言，检测框边界的细微变化并不一定会直接影响系统运行，目标中心点的连续性才是控制模块更关注的输入信息。因此，检测模块输出结果需要与后续控制策略配合使用，不能简单地把单帧检测框看作最终跟踪结果。

\subsection{ReID 目标记忆测试}

ReID 模块包含独立单元测试。\texttt{tests/test\_body\_reid.py} 验证了人体特征归一化、短时离开后合并为同一 ID、不同人体分配新 ID、同一帧中不重复分配同一历史 ID、指定 ID 存在时优先选择以及指定 ID 缺失时等待等逻辑。\texttt{tests/test\_body\_tracking.py} 验证了目标选择函数能够选择面积最大的人体，并丢弃过期检测元数据。

这些测试覆盖了 ReID 记忆模块的核心逻辑，可以看出目标 ID 分配和指定目标选择在纯逻辑层面具有可验证性。该模块采用轻量颜色和几何特征，适合在低开销条件下支持短时目标记忆；对于遮挡、服装颜色相近和光照剧烈变化等更复杂场景，也为后续结合外观特征或运动模型提供了接口基础。

单元测试的意义在于将目标记忆逻辑从真实设备环境中抽离出来。即使没有连接布控球，也可以通过构造检测框、面积、中心点和特征向量来验证 ID 分配是否符合预期。这种测试方式能够降低调试成本，也能在后续修改 ReID 模块时，帮助确认已有核心行为没有被破坏。

从测试结果看，轻量 ReID 记忆模块能够为原型系统提供基本的目标连续性支持。它适合处理目标短时遮挡、短时间离开画面后重新出现等情况，在多人场景下能够为目标选择层提供轻量增强。本文将该模块定位为自动跟踪系统中的身份辅助单元，使系统在保持实时性的同时获得更稳定的目标编号。

\subsection{PTZ 自动跟踪测试}

PTZ 测试用于验证目标偏差能否转换为云台运动指令。测试过程中，使目标在画面中产生水平或垂直方向位移，并观察目标偏移量变化与云台动作方向是否一致。系统通过阈值判断减少中心附近的小幅抖动，通过比例脉冲使偏差较大时转动时间略长，再通过冷却时间限制指令频率。根据中期检查结果，系统已经实现基本 PTZ 跟踪功能，目标位于画面边缘时，云台能够朝目标方向调整。

该测试重点关注三个现象：一是当目标接近画面中心时，云台应尽量保持静止，避免因检测框轻微变化产生频繁抖动；二是当目标偏离中心较远时，云台应能够朝正确方向动作，使目标逐步回到画面中心附近；三是当用户关闭 PTZ 自动跟踪或切换到手动控制时，系统应立即停止自动控制逻辑，避免自动指令与人工操作冲突。

在实际观察中，PTZ 自动跟踪效果与检测稳定性和云台机械响应都有关系。即使目标检测输出正确，如果控制指令过密，云台仍可能出现不必要的启停；如果控制指令过少，目标回到画面中心的速度又会变慢。因此，PTZ 测试不仅要看方向是否正确，还要观察系统在目标接近中心时是否能够及时停止，以及在目标偏离较大时是否能够保持足够的响应速度。

\subsection{交互模块测试}

交互模块测试包括检测状态切换、PTZ 控制状态切换、人工接管、目标选择、检测结果反馈和倒挂模式切换。测试结果显示，系统能够通过交互模块控制 YOLO 检测和自动跟踪状态，也能够在必要时进行人工接管。目标选择逻辑能够在自动模式和可见目标 ID 之间切换，指定 ID 不可见时，系统会等待该目标重新出现。倒挂模式能够翻转预览画面和检测画面，并反转控制偏差，从而适配布控球倒悬安装场景。

从使用体验看，交互模块测试反映出系统不仅具备算法功能，也具备基本调试可用性。检测结果反馈能够展示目标框和偏移量，帮助判断控制输入是否合理；状态提示能够反映当前检测与跟踪模式，降低误操作概率；人工接管能力能够为现场调试提供补充控制方式。这些设计有助于系统在真实设备调试中更快确认运行状态。

交互模块还可以帮助分析系统内部状态是否一致。例如，当界面显示的目标 ID 与控制线程选中的目标一致时，说明检测结果、目标记忆和控制决策之间的数据传递正常；当倒挂模式开启后，若画面方向和云台动作方向都发生对应变化，则说明倒挂适配逻辑生效。通过这些可视化反馈，测试人员可以在不阅读底层日志的情况下快速判断系统是否处于预期状态。

\section{性能分析}

从中期检查和当前实现看，系统在实时性方面已经达到原型验证要求。视频流链路能够稳定获取实时画面，检测模块使用 YOLOv8n 轻量模型，推理时间约为 30 ms，检测帧率可达到 30 fps 以上。PTZ 控制线程以约 50 Hz 的频率检查检测结果，并通过 0.012 s 的冷却时间限制控制指令发送频率，避免连续发送造成设备抖动。

需要说明的是，实时性并不只由模型推理时间决定。布控球视频采集、播放库解码、帧格式转换、检测结果更新、控制线程调度和设备机械响应，都会影响用户最终感受到的跟踪效果。因此，本文采用“读取最新帧”和“控制冷却”的思路，在保证系统响应速度的同时尽量减少积压和抖动。

影响系统端到端响应的主要因素包括视频解码延迟、图像文件读写开销、YOLO 推理耗时、线程调度开销和布控球机械响应时间。当前系统通过读取最新帧而不是逐帧排队处理，降低了检测线程因积压帧造成的延迟；在进一步工程优化时，也可以将帧数据直接在内存中转换为 OpenCV 图像，以继续提升链路效率。

从控制稳定性看，系统当前采用阈值死区、指数平滑和短脉冲控制，在常规移动速度下能够缓解大部分抖动。对于目标快速穿过画面、检测框短时跳变或云台机械响应较慢等场景，系统可以继续通过运动预测和控制参数自适应来提升跟踪平滑性。由此可见，后续工程优化可以从检测模型、目标运动估计和控制策略三个方面协同展开。

从资源占用角度看，当前系统保留了便于调试的中间帧保存方式，例如将解码图像保存为中间文件后再读取。这种方式有利于开发阶段检查帧内容和定位链路状态。在嵌入式平台迁移场景下，可以进一步改为内存帧传递，并根据实际算力调整检测频率和输入尺寸，以减少端到端延迟。

\section{问题分析与优化方向}

结合上述测试结果，系统工程优化可以从以下几个方面展开：

\begin{enumerate}
    \item 跟踪平滑性优化。对于快速移动目标或检测框波动较大的场景，可以在现有指数平滑和阈值死区基础上，引入滑动平均、Kalman 滤波或速度变化率限制，使云台动作更加连续。
    \item 复杂场景适应性增强。针对遮挡、逆光、强光变化和多人密集场景，可以加入亮度均衡、去噪、数据增强训练和更鲁棒的检测模型，使人体检测和 ReID 特征匹配在多种场景下保持较好的稳定性。
    \item 系统资源调度优化。YOLO 推理和图像读写会占用一定 CPU/GPU 资源，可以通过 INT8 量化、TensorRT 加速、检测频率动态调节和多线程调度优化，提高端到端运行效率。
    \item 目标身份特征扩展。当前 ReID 模块使用颜色和几何特征，适合短时记忆。进一步扩展时，可以接入向量数据库和深度人体特征提取模型，记录人物出现时间并形成更完整的目标记忆机制。
\end{enumerate}

除上述方向外，系统还可以继续开展长时间运行验证。当前测试主要集中在功能闭环和短时间运行效果，进一步覆盖长时间连续预览、网络波动、设备重连、模型异常和系统资源释放等情况，有助于更完整地评价系统运行状态。若系统部署到真实巡检或安防场景，可以增加日志记录、异常告警、自动重连和运行状态监控等机制。

测试指标记录方式也可以进一步细化。例如，在视频链路中记录帧获取时间和检测开始时间，在检测模块中记录单帧推理耗时和有效检测数量，在控制模块中记录指令发送时间和目标偏差变化。通过这些日志，可以更清楚地分析端到端延迟来源，也便于比较不同模型、不同输入尺寸和不同控制参数下的系统表现。

\section{测试结论}

综合功能测试和初步性能分析，本文系统已经实现布控球视频接入、YOLO 人体检测、人体 ID 记忆、自动/指定目标跟踪、PTZ 控制和图形界面集成，形成了完整的智能跟踪原型。系统能够满足基本实时检测与云台跟踪需求，并具备向复杂场景适应、资源调度优化和长期运行管理继续扩展的基础。

总体而言，测试结果验证了本文系统方案的可行性。视频采集层能够稳定提供输入，检测分析层能够输出人体目标和 ID，控制决策层能够根据偏差生成云台动作，用户界面也能够支持状态切换和人工干预。在此基础上，系统可以继续围绕模型优化、控制参数自适应和嵌入式部署开展工程拓展。
